\section{Related Works}

\paragraph{Listening head generation.}
Listening head generation synthesizes non-verbal listener behaviors conditioned on speaker-side conversational cues. ViCo~\cite{zhou2022responsive} first established responsive listening head generation as a benchmark, and Learning2Listen~\cite{ng2022learning} modeled listener motion as a non-deterministic distribution to reflect the multi-valid nature of dyadic interaction. Later methods improved listener synthesis with emotion conditioning~\cite{song2023emotional}, non-autoregressive sequence modeling~\cite{liu2024listenformer}, diffusion-based generation~\cite{tran2024dim}, dyadic interaction modeling~\cite{wang2025diffusion}, and real-time multimodal interaction~\cite{zhu2025infp,guo2025arig}. However, most existing methods optimize holistic motion realism, diversity, or speaker-listener synchrony, without explicitly modeling listener reactions as categorized and temporally localized events. Our work instead focuses on generating reactions with appropriate type, timing and duration.

\paragraph{Interactive conversation datasets.}
Dyadic conversation datasets provide important resources for modeling social interaction. Early corpora such as RECOLA~\cite{ringeval2013recola} and NoXi~\cite{cafaro2017noxi} support affective and feedback behavior analysis, while ViCo~\cite{zhou2022responsive} introduced a task-specific benchmark for responsive listening heads. Larger resources such as RealTalk~\cite{geng2023realtalk}, Seamless~\cite{agrawal2025seamless}, and SpeakerVid-5M~\cite{zhang2025speakervid5m} provide more diverse conversational videos, but they are not designed around fine-grained listener reaction events. In contrast, our dataset constructs aligned speaker-listener clips with explicit event-level annotations of reaction type, timing, and duration.

\paragraph{Evaluation methods.}
Listener generation is commonly evaluated using backchannel prediction metrics, such as precision, recall, F1, tolerance windows, and subjective appropriateness judgments~\cite{morency2008predicting,dekok2012survey}, or visual generation metrics such as FID~\cite{Seitzer2020FID}, FVD~\cite{unterthiner2018fvd}, and LPIPS~\cite{zhang2018lpips}. These metrics measure feedback timing or global visual quality, but do not jointly assess whether a generated listener reacts with the correct type, at the correct moment, and with realistic reaction-specific motion. Recent reaction generation benchmarks emphasize appropriateness, diversity, synchrony, and realism~\cite{song2023fmarg,song2024react}. Our evaluation further treats reactions as categorized temporal events and measures reaction occurrence, temporal alignment, asymmetric timing deviation, and reaction-region quality.